\section{Introduction}
\label{sec:intro}

Speech has become a primary interface to cloud intelligence. Voice assistants,
meeting transcription products, live captioning services, and multilingual
communication platforms are routinely assembled from commercial automatic
speech recognition (ASR) and machine translation (MT) interfaces offered by
large cloud providers. In a typical deployment, audio captured by a far-field
microphone or transported over a communication channel is first converted to
text by an ASR service, and the resulting text is then passed to a text
translation engine whenever the application crosses a language boundary. This
cascade architecture dominates industrial practice because it allows each
stage to be selected, replaced, and billed independently, and because the best
available components at each stage are often operated by different providers.
An alternative architecture, end-to-end speech translation (E2E ST), maps
source-language audio directly to target-language text and has matured
rapidly with large foundation models.

The acoustic pathway between a speaker and these services is rarely pristine.
Speech arrives mixed with competing talkers, shaped by room reverberation,
band-limited by telephone channels, compressed by low-bitrate codecs, or
corrupted by packet loss on unstable network links. The robustness of ASR
systems to such degradations has been a central concern of the speech research
community for decades, and controlled robustness corpora have driven
substantial progress. What has changed is the identity of the systems that
now dominate deployment. Commercial speech services are increasingly built on
large language model (LLM) based ASR, on general-purpose multimodal foundation
models that accept audio alongside text, and on streaming speech translation
models, none of which can be fine-tuned, inspected, or benchmarked offline by
their consumers. A systems integrator who composes such services into a
translation pipeline has no principled way to predict how recognition errors
induced by an acoustic channel will propagate into the final translation, or
whether a cascade or an end-to-end design is the safer choice under the
conditions of a given deployment.

The evaluation gap is not academic. Deployment decisions are made
weekly by integrators who must choose among service tiers and architectural
patterns with nothing but clean-audio demo pages and marketing claims, while
the operating conditions of real audio pathways, conference rooms, vehicles,
telephone trunks, and congested wireless links, are precisely the conditions
that clean demos do not exercise. A deployment that selects a recognition
frontend on clean-audio accuracy alone can be quietly losing most of its
translation quality in the field, and a deployment that assumes end-to-end
robustness can be more fragile than the cascade it replaced. Closing this
gap requires measurement that is controlled at the acoustic layer, complete
at the pipeline layer, and reproducible at the service layer, and that is
the design brief of the benchmark presented here.

This gap motivates the present study. Academic robustness benchmarks such as
Aurora~\cite{pearce2000aurora} and the CHiME challenge
series~\cite{barker2018fifth} evaluate open recognizers on fixed noise
corpora, but they do not cover the black-box commercial systems that now carry
production traffic, they rarely extend beyond recognition to the downstream
translation stage, and they cannot be re-run on new model versions without
access to model internals. Conversely, studies that compare commercial ASR
services report error rates on fixed test sets but typically consider clean
audio and stop at the transcription boundary. The propagation of acoustic
degradation through a complete speech translation pipeline, comparing
architectural families and both cascade and end-to-end designs under
identical and fully controlled channel conditions, has to our knowledge not
been quantified for the current generation of speech APIs.

This paper presents CascadeBench, a black-box benchmark designed to answer
four questions. First, how sensitive are the main product classes of
commercial ASR, namely specialised streaming models, LLM-based recognisers,
and general-purpose multimodal models, to eight canonical acoustic channel
degradations? Second, how does recognition accuracy scale with noise
severity across the full range from mild to destructive, and where does
each class fail? Third, when
recognition errors flow into a downstream translation engine, does the
cascade amplify them, and how does the cascade compare with end-to-end speech
translation under identical conditions? Fourth, what are the latency
consequences of each choice? To answer these questions under the
constraint that only public interfaces are available, the benchmark
synthesises a stratified bilingual corpus of 480 Mandarin Chinese and English
sentences with multiple neural voices, applies a grid of eight controlled
acoustic degradations together with an eight-level noise severity sweep from
20 to $-10$~dB, and measures five recognition systems, three cascade
translation paths, and two
end-to-end speech translation paths through their public APIs, using only
free service quotas. All systems come from one provider, so the study is a
single-provider case study whose class-level findings require cross-provider
replication before attribution to architecture; this constraint is examined
in Section~\ref{sec:limitations} and revisited throughout.

The study yields findings that are relevant to both the science of robust
speech processing and the practice of assembling speech services. Recognition
services from the three classes differ sharply in where, rather than how
fast, they fail under noise: the LLM-based recognisers hold a sub-percent
error floor and cross ten percent error at noise levels 2.5 to 4~dB lower
than the specialised streaming system, with sharper collapse beyond their
knee, and below 0~dB they fail by insertion-heavy hallucination that drives
character error rates above 100 percent. The end-to-end translation paths,
often assumed to be more robust
than cascades because they avoid error propagation, do not inherit
robustness automatically: the omni-modal path loses more of its clean
quality under severe noise than a well-frontend cascade, while a streaming
end-to-end service is the most resilient path measured. Across fifteen
operating conditions, the BLEU score of every cascade measured
falls linearly with its recognition error rate, with a pass-through slope of
0.44 to 0.54 BLEU points per percentage point, sentence-level bootstrap
confidence intervals within $\pm 0.04$, and an intercept statistically
indistinguishable from zero, which gives integrators a calibrated
engineering rule for estimating translation losses from recognition
quality alone. Finally, the benchmark surfaces behavioural differences that
are invisible in clean-audio evaluations, including two distinct
empty-output mechanisms under severe noise, an explicit protocol-level
refusal and a silent empty return, and the unique fragility of a
telephone-band specialist under reverberation.

The contribution of this paper is threefold. It provides, to our knowledge,
the first systematic measurement of acoustic degradation propagation across
recognition and translation stages for commercial black-box speech APIs,
covering three product classes of one provider and both cascade and
end-to-end translation designs, which we frame as a case study pending
cross-provider replication. It introduces a fully reproducible benchmark
recipe in which the test material, the degradations, and every model response
are controlled, logged, and released, so that the measurements can be
repeated as providers update their models. And it distils the measurements
into quantitative design guidance, including noise-failure thresholds per
product class, the error pass-through regression for cascades, and a
dual-reference protocol that exposes and controls the systematic bias that
arises when automatic translation quality is scored against references
produced by models related to one of the systems under test.

The remainder of the paper is organised as follows.
Section~\ref{sec:related} reviews related work on robust ASR, cascade and
end-to-end speech translation, and the evaluation of commercial services.
Section~\ref{sec:methods} describes the benchmark design, including the
corpus, the degradation grid, the systems under test, and the evaluation
protocol. Section~\ref{sec:results} presents the results, and
Section~\ref{sec:discussion} discusses their implications for practitioners
and for future evaluation practice. Limitations are examined in
Section~\ref{sec:limitations}, and Section~\ref{sec:conclusions} concludes.
